\documentclass[10pt]{article}
\usepackage[T1]{fontenc}\usepackage[utf8]{inputenc}\usepackage{mathptmx}
\usepackage[margin=0.8in]{geometry}
\setlength{\parindent}{0pt}
\pagestyle{empty}
\begin{document}
\begin{center}
{\large\bfseries Weighted PageRank and Official Rankings in Tennis and Other Sports}\\[6pt]
Chakib Madcur\\[2pt]
{\small Department of Computer Science, Roosevelt University, Chicago, IL, USA}
\end{center}
\vspace{0.6em}
\textbf{Abstract.} Official sports rankings follow different rules: tournament points in tennis, an Elo-type rating in international football, and league tables in team sports. Using ideas from network science, I rank competitors with weighted PageRank on the directed network of match results, where every loser points to the winner and each link is weighted by the number of wins. I compare these PageRank rankings with the official rankings over five decades of men's and women's professional tennis, and in international football, the Premier League, the NBA, the NHL and the NFL. The two rankings are well correlated, so I focus on their differences. I compare their predictability by using each ranking at the end of one season to predict the winners of the next season's matches, and I examine where the rankings disagree: how much one upset of the world number one is worth, and which players gain or lose places depending on where and in which tournaments they play. The aim is to understand what each ranking rule rewards.
\end{document}
